% Part: history
% Chapter: biographies
% Section: ernst-zermelo

\documentclass[../../../include/open-logic-section]{subfiles}

\begin{document}

\olfileid{his}{bio}{zer}

\olsection{恩斯特·策梅洛}

\olphoto{zermelo-ernst}{恩斯特·策梅洛}

恩斯特·策梅洛于1871年7月27日出生在德国柏林。他有五个姐妹，但家人健康状况不佳，仅有三人活至成年。他的父母在他年幼时相继去世，致使他十七岁时便与兄弟姐妹沦为孤儿。策梅洛对艺术，尤其是诗歌，有着浓厚的兴趣。他以敏锐、机智和富有批判精神而闻名。他最为人称道的数学成就包括选择公理的引入（1904年）以及集合论的公理化（1908年）。

策梅洛在大学时期的兴趣十分广泛。他修读了物理、数学和哲学课程。在赫尔曼·施瓦茨的指导下，策梅洛于1894年在柏林大学完成了他的博士论文《变分法研究》。1897年，他决定前往哥廷根大学继续深造，在那里他深受大卫·希尔伯特基础研究工作的影响。1899年他取得了教授资格，但直到十一年后才获得教授职位——这可能是因为他古怪的举止和“神经质的急躁”。

策梅洛最终于1910年在苏黎世大学获得了带薪教授职位，但因患肺结核于1916年被迫退休。康复后，他于1921年在弗赖堡大学获得荣誉教授职位。在此期间，他致力于数学基础的研究。他对托拉尔夫·斯科伦和库尔特·哥德尔的工作感到恼火，并在自己的论文中公开批评他们的方法。由于不受欢迎以及反对希特勒在德国掌权，他于1935年被弗赖堡大学解职。

策梅洛的晚年十分孤寂。1935年被解职后，他放弃了数学。他搬到乡下，过着简朴的生活。他于1944年结婚，并因视力逐渐衰退而完全依赖妻子照料。到1951年，策梅洛完全失明。他于1953年5月21日在德国金特斯塔尔（Günterstal）去世。

\begin{reading}
关于策梅洛的完整传记，请参阅 \citet{Ebbinghaus2015}。
策梅洛1904年与1908年开创性论文的德文原文可供阅读 \citep{Zermelo1904,Zermelo1908}。策梅洛的文集（包括他关于物理学的著述）有英译本可供阅读 \citep{Ebbinghaus2010,Ebbinghaus2013}。
\end{reading}

\end{document}